% ============================================================
% Section 178: 类氢原子n=2能级在电磁场中的分裂
% ============================================================
\section{量子力学：类氢原子\texorpdfstring{$n=2$}{n=2}能级在电磁场中的分裂}
\label{sec:hydrogen-n2-em-splitting}

\begin{modelbox}[题目：类氢原子\texorpdfstring{$n=2$}{n=2}能级在电磁场中的分裂]
考虑类氢原子（核电荷$Ze$，略去自旋和相对论效应）的$n=2$能级。
\begin{enumerate}
    \item 计算存在均匀磁场时能级分裂的最低阶；
    \item 存在均匀电场时做类似计算；
    \item 电场、磁场同时存在并相互垂直时，做同样的计算（径向波函数和角向积分可用参数代替）。
\end{enumerate}
\end{modelbox}

\begin{tipbox}[考点快速分析]
\begin{itemize}
    \item \textbf{塞曼效应}：$H'_B=\dfrac{eB}{2m_ec}L_z$，选择定则$\Delta m=0$，正常塞曼分裂
    \item \textbf{Stark效应}：$H'_E=e\mathcal{E}z$，选择定则$\Delta l=\pm1$、$\Delta m=0$
    \item \textbf{正交电磁场}：$H'=H'_B+H'_E$，联合对角化
    \item \textbf{参数法}：用$\beta$、$\gamma$代替具体积分值
\end{itemize}
\end{tipbox}

\begin{derivationbox}[标准解题过程]
\textbf{1. 均匀磁场（正常塞曼效应）}

取 $\mathbf{B}=B\hat{\mathbf{z}}$，微扰 $H'_B=\dfrac{eB}{2m_ec}L_z$。在 $|nlm\rangle$ 基中 $L_z$ 对角，本征值 $m\hbar$。

$n=2$ 的四个态：
\begin{itemize}
    \item $|200\rangle$、$|210\rangle$：$m=0$，修正为零；
    \item $|211\rangle$：$m=1$，$E^{(1)}=+\dfrac{eB\hbar}{2m_ec}\equiv+\beta$；
    \item $|21,-1\rangle$：$m=-1$，$E^{(1)}=-\beta$。
\end{itemize}

\begin{equation}
\boxed{n=2\text{ 分裂为三条：}E_2^{(0)},\;E_2^{(0)}\pm\frac{eB\hbar}{2m_ec}.}
\end{equation}

\textbf{2. 均匀电场（线性Stark效应）}

取 $\boldsymbol{\mathcal{E}}=\mathcal{E}\hat{\mathbf{z}}$，$H'_E=e\mathcal{E}z=e\mathcal{E}r\cos\theta$。非零矩阵元仅 $\langle200|H'_E|210\rangle=-\lambda$，其中$\lambda=3e\mathcal{E}a_0/Z$。

微扰矩阵对角化得：
\begin{equation}
\boxed{E^{(1)}=+\frac{3e\mathcal{E}a_0}{Z},\;-\frac{3e\mathcal{E}a_0}{Z},\;0,\;0.}
\end{equation}

\textbf{3. 正交电磁场}

设 $\mathbf{B}=B\hat{\mathbf{z}}$，$\boldsymbol{\mathcal{E}}=\mathcal{E}\hat{\mathbf{x}}$。$x=r\sin\theta\cos\varphi=\frac{r}{2}\sin\theta(e^{i\varphi}+e^{-i\varphi})$。$H'_E$ 的非零矩阵元满足 $\Delta l=\pm1$、$\Delta m=\pm1$。

定义：
\[
\beta=\frac{eB\hbar}{2m_ec},\qquad\gamma=\frac{3e\mathcal{E}a_0}{\sqrt2\,Z}.
\]

在基 $\{|211\rangle,|21,-1\rangle,|210\rangle,|200\rangle\}$ 中，总微扰矩阵：
\[
H'=\begin{pmatrix}\beta&0&0&-\gamma\\0&-\beta&0&\gamma\\0&0&0&0\\-\gamma&\gamma&0&0\end{pmatrix}.
\]

$|210\rangle$ 完全解耦（能量修正为零）。剩余 $3\times3$ 久期方程给出
\[
E^{(1)}\bigl[(E^{(1)})^2-(\beta^2+2\gamma^2)\bigr]=0.
\]

解得：
\begin{equation}
\boxed{E^{(1)}=0,\;\pm\sqrt{\beta^2+2\gamma^2}=\pm\sqrt{\Bigl(\frac{eB\hbar}{2m_ec}\Bigr)^2+\Bigl(\frac{3e\mathcal{E}a_0}{Z}\Bigr)^2}.}
\end{equation}

加上 $|210\rangle$ 的零修正，$n=2$ 分裂为三条。当 $B=0$ 退化为Stark分裂；当 $\mathcal{E}=0$ 退化为Zeeman分裂。
\end{derivationbox}

\begin{formulabox}[答案汇总]
\begin{center}
\begin{tabular}{ll}
\toprule
\textbf{结论} & \textbf{结果} \\
\midrule
纯磁场 & $E^{(1)}=0,\;\pm eB\hbar/(2m_ec)$ \\
纯电场 & $E^{(1)}=0,\;\pm3e\mathcal{E}a_0/Z$ \\
正交电磁场 & $E^{(1)}=0,\;\pm\sqrt{\beta^2+2\gamma^2}$ \\
\bottomrule
\end{tabular}
\end{center}
\end{formulabox}

\begin{insightbox}[物理图像]
\begin{itemize}
    \item 正交电磁场下的联合分裂公式是塞曼效应和Stark效应的``几何叠加''，类似于两个垂直振子的频率合成。这体现了微扰论中不同微扰的可加性。
    \item 该模型是理解交叉场（crossed-field）光谱和量子Hall效应的前身。在交叉场中，电子运动形成摆线（cycloid）轨迹，量子化后产生复杂的能带结构。
    \item $1/Z$ 的标度律说明类氢离子（如He$^+$、Li$^{2+}$）的Stark分裂比氢原子小，因为更高核电荷使电子更紧密束缚，更难被外场极化。
\end{itemize}
\end{insightbox}
